\scanpage{109}
\mspage{A 3}
% archive sheet number: 108
% Faint reversed diagram strokes in the upper margin are show-through.
ou moins grossières, plus ou moins subtiles, suivant l'intéressé. Je ne prétends nullement y faire exception. La lecture de mon témoignage ne laissera, j'espère, aucun doute à ce sujet.

Il est vrai aussi que l'ambition la plus dévorante est impuissante à découvrir le moindre énoncé mathématique, ou à le démontrer - tout comme elle est impuissante (par exemple) à ``faire bander'' (au sens propre du terme). Qu'on soit femme ou homme, ce qui ``fait bander'' n'est nullement l'ambition, le désir de briller, d'exhiber une puissance, sexuelle en l'occurence - bien au contraire ! Mais c'est la perception aigüë de quelque chose de fort, de très réel et de très délicat à la fois. On peut l'appeler ``la beauté'', et c'est là un des mille visages de cette chose-là. D'être ambitieux n'empêche pas forcément de sentir parfois la beauté d'un être, ou d'une chose, d'accord. Mais ce qui est sûr, c'est que ce n'est \mbox{p a s} l'ambition qui nous la fait sentir...

L'homme qui, le premier, a découvert et maîtrisé le feu, était quelqu'un exactement comme toi et moi. Pas du tout ce qu'on se figure sous le nom de ``héros'', de ``demi-dieu'' et j'en passe. Sûrement, comme toi et comme moi, il a connu la morsure de l'angoisse, et la pommade vaniteuse éprouvée, qui fait oublier la morsure. Mais au moment où il a ``connu'' le feu, il n'y avait ni peur, ni vanité. Telle est la vérité dans le mythe héroïque. Le mythe devient insipide, il devient pommade, quand il nous sert à nous cacher un \mbox{a u t r e} aspect des choses, tout aussi réel et tout aussi essentiel.

Mon propos dans Récoltes et Semailles a été de parler de l'un et de l'autre aspect - de la pulsion de connaissance, et de la peur et de ses antidotes vaniteux. Je crois ``comprendre'', ou du moins \mbox{c o n n a î t r e} la pulsion et sa nature. (Peut-être un jour découvrirai-je, émerveillé, à quel point je me faisais illusion ...) Mais pour ce qui est de la peur et de la vanité, et les insidieux blocages de la créativité qui en dérivent, je sais bien que je n'ai pas été au fond de cette grande énigme. Et j'ignore si je verrai jamais le fond de ce mystère, pendant les années qui me restent à vivre ...

En cours d'écriture de Récoltes et Semailles deux images ont émergé, pour représenter l'un et l'autre de ces deux aspects de l'aventure humaine. Ce sont \mbox{l ' e n f a n t} (alias \mbox{l' o u v r i e r}), et le \mbox{P a t r o n}. Dans la Promenade qu'on va faire tantôt, c'est de ``l'enfant'' qu'il sera question presque exclusivement. C'est lui aussi qui figure dans le sous-titre ``\mbox{L ' e n f a n t \quad e t \quad l a \quad M è r e}''. Ce nom va s'éclairer, j'espère, au cours de la promenade.
